%======================================================================
% CiteREAL: Retrieval-Augmented Citation Intent Classification
% -----------------------------------------------
% arXiv-ready manuscript, CVPR-style two-column layout.
%
% HOW TO COMPILE
%   Option A (official CVPR look): open Overleaf, start from the official
%   "CVPR template" project, replace main.tex with this file, and set
%   \usecvprtrue below. Compile with pdflatex (run twice).
%   Option B (self-contained): keep \usecvprfalse; the file defines a
%   compatible title/author/abstract block and compiles standalone.
%
% BEFORE ARXIV UPLOAD
%   1. Replace the code-availability wording with your real repository URL
%      (search for "TODO-REPO" below).
%   2. All numbers are real measured values from runs/*/metrics.json and
%      runs/fast_bm25l_bench.json.
%======================================================================
\documentclass[10pt,twocolumn,twoside]{article}

% ---- switch: set 1 to load the official CVPR style, 0 for plain ----
\newif\ifusecvpr
\usecvprfalse   % <-- set \usecvprtrue if you have cvpr.sty (official kit)

\ifusecvpr
\usepackage{cvpr}                    % camera-ready; use [review] for blind
\bibliographystyle{ieee_fullname}
\else
% ---------- self-contained CVPR-like two-column layout ----------
\usepackage[margin=1in]{geometry}
\usepackage[T1]{fontenc}
\usepackage[utf8]{inputenc}
\usepackage{times}                   % Times for the CVPR look
\usepackage{amsmath,amssymb,amsfonts}
\usepackage{graphicx}
\usepackage{booktabs}
\usepackage{multirow}
\usepackage{array}
\usepackage{textcomp}
\usepackage{caption}
\captionsetup{font=small,labelfont=bf}
\usepackage[colorlinks=true,urlcolor=blue,citecolor=blue]{hyperref}
\setlength{\columnsep}{0.55cm}
% CVPR-style title/author/abstract block
\makeatletter
\def\maketitle{%
\twocolumn[{%
\thispagestyle{empty}%
\begin{center}%
  \let\footnote\thanks
  {\LARGE \bfseries \@title\par}%
  \vskip 0.4em%
  {\large\lineskip .5em%
   \begin{tabular}[t]{c}\@author\end{tabular}\par}%
  \vskip 0.5em%
\end{center}%
\vskip 1em%
}]%
}%
\makeatother
\fi

\title{When Retrieval Does Not Help: A Controlled Study of\\
       Citation Intent Classification with BM25L Evidence}

\author{
Chen Boxin\\
School of Business Administration, Huaqiao University\\
Quanzhou 362000, China\\
\texttt{2416411001@stu.hqu.edu.cn}
}

\begin{document}
\maketitle

%======================================================================
\begin{abstract}
Citation intent classification labels the sentence containing a citation
marker as \emph{background}, \emph{method}, or \emph{result} use. We test a
natural idea on this task: prepend the input with $k$ similar contexts
retrieved by BM25L, so the encoder sees the target sentence together with
lexical evidence from neighbours in the training pool. Across BERT, SciBERT and
DeBERTa-v3-base on SciCite, under a leakage-controlled protocol, the answer is
negative. DeBERTa-v3-base alone reaches 86.89\% accuracy and 85.49\% macro-F1,
and no augmented variant beats it. Three results qualify that headline. First,
retrieval-augmented inputs break a standard \mbox{[CLS]} head (accuracy falls
to 59.97\%), while a gated dual-pooling head that fuses the sentence vector
with a pooled evidence vector restores 86.14\% / 84.71\%, in effect recovering
the baseline. The head, not the evidence, is the fragile part. Second, the
method is cheap once the engineering is done: a sparse-matrix rewrite of the
BM25L scorer cuts a full preprocessing pass over 8{,}243 contexts from about
8 minutes to about 12 seconds, cached re-builds take under one second, and
training throughput is unchanged on a single 6\,GB GPU. Third, the one
configuration that appears to beat the backbone---a mask that forbids
same-label evidence, lifting accuracy to 96.40\%---does so because the mask
itself leaks: with three classes and $k{=}3$ neighbours, the target label is
exactly the class absent from the evidence, and the model learns the
elimination rule. All code,
configurations and run logs are released.
\end{abstract}

%======================================================================
\section{Introduction}\label{sec:intro}

Citation intent classification (CIC) asks a question that bibliometric tools
increasingly need answered automatically: given the sentence in which a
citation appears, why is the cited work referenced? Cohan et~al.~\cite{cohan2019}
framed the task over three intents (\emph{background}, \emph{method},
\emph{result}) and released SciCite, 11k annotated contexts spanning computer
science and biomedicine, which remains the standard benchmark.

Pretrained encoders dominate current leaderboards. SciBERT~\cite{beltagy2019}
adapted BERT to scientific text; DeBERTa-v3~\cite{he2021} replaced absolute
position embeddings with disentangled attention and transfers well across
classification benchmarks. On the retrieval side, models such as
RAG~\cite{lewis2020} and DPR~\cite{karpukhin2020} showed that an external
corpus, consulted at inference time, lifts performance on open-domain
question answering, precisely where a single input lacks the information
needed to answer.

CIC sits between the two. The input is one sentence, and intent is often
recoverable from surface cues: ``we apply the method of X'' reads as a method
citation, ``in contrast to Y'' as a result comparison. But some contexts are
ambiguous in isolation, and similar contexts elsewhere in the corpus carry the
same phrases with the label attached. Retrieval augmentation therefore has a
plausible mechanism here, and the mechanism is cheap to test: retrieve
neighbours with BM25, concatenate, fine-tune.

The mechanism, however, has an obvious failure path: if the retrieved
neighbour's label equals the target's, the model can copy rather than learn.
Leakage of this kind is easy to introduce and easy to miss. So we ran the
study with the safeguards built in (train-only pool, $O(1)$ self-exclusion,
and a leak-free mask that forbids same-label evidence) and with three
backbones of increasing strength.

The outcome is a clean negative. With safeguards in place, no
retrieval-augmented variant beats its plain encoder, and the strongest
backbone, DeBERTa-v3-base, stays ahead at 86.89\% / 85.49\%. The study
produced two findings worth reporting in their own right. The first is
diagnostic: concatenating heterogeneous evidence silently breaks a standard
\mbox{[CLS]} head (59.97\% accuracy), and a gated dual-pooling head fixes this,
restoring the augmented model to within 0.8 points of its backbone. If one
runs retrieval augmentation and sees it fail, or sees it ``work'' once a
leakage control is removed, the pooling head and the retrieval pool are the
first places to look. The second is practical: the full pipeline, including
a fresh preprocessing pass over all 8{,}243 training contexts, runs in
seconds on the same 6\,GB GPU that trains the models, so reruns and ablations
cost almost nothing once the caching layer is in place.

Our contributions:
\begin{itemize}
  \item a leakage-controlled evaluation of BM25L retrieval augmentation for
        CIC over three encoders, with a negative result against a strong
        DeBERTa-v3 baseline;
  \item a gated dual-pooling head that stabilises retrieval-augmented inputs,
        with an ablation showing the magnitude of the CLS-head failure it
        prevents;
  \item an engineering account with measured numbers: a vectorised BM25L
        scorer (76$\times$ faster at scoring and 43$\times$ faster at building
        the augmented corpus than the reference implementation), a disk cache
        for augmented inputs, and unchanged training throughput on 6\,GB of
        VRAM.
\end{itemize}

%======================================================================
\section{Related Work}\label{sec:related}

\subsection{Citation function analysis}
Teufel et~al.~\cite{teufel2006} classified citation function with
argumentative zones and hand-crafted features. Jurgens et~al.~\cite{jurgens2018}
built ACL-ARC, six classes over computational linguistics papers, and traced
how a field's citation practice evolves. SciCite~\cite{cohan2019} scaled the
annotation effort to 11k contexts and added structural scaffolds (section
name, citing-cited paper pairs) on top of ELMo~\cite{peters2018}. Domain
transfer and class imbalance, not model capacity, turned out to be the hard
parts, and subsequent work has largely kept the dataset as the reference
benchmark.

\subsection{Pretrained encoders for scientific text}
BERT~\cite{devlin2019} set the template; SciBERT~\cite{beltagy2019} pre-trained
from scratch on Semantic Scholar corpora; DeBERTa-v3~\cite{he2021} combined
disentangled attention with ELECTRA-style pre-training~\cite{clark2020} and
replaced the absolute position bias with relative ones. For sentence-level
classification tasks the encoder choice now accounts for several points of
macro-F1, a gap that any augmentation method has to clear before it counts as
a contribution. Graph-structured and citation-network
features~\cite{jurgens2018} help when linkage data exists, but they are
unavailable in the settings we target, where only the citing context is
reliably present.

\subsection{Retrieval augmentation}
REALM~\cite{guu2020} and RAG~\cite{lewis2020} retrieve-then-read for
knowledge-intensive generation; DPR~\cite{karpukhin2020} learns the retriever
itself. Outside QA, sparse retrievers still find use because their evidence is
cheap and inspectable: BM25~\cite{robertson2009} and its length-normalised
variant BM25L~\cite{lv2011} remain strong baselines. Our study is closer in
spirit to controlled studies of retrieval's benefit: we hold the backbone
fixed, vary the retrieval pool under leakage controls, and measure what
changes. We are not aware of a prior CIC study that reports the interaction
between evidence concatenation and pooling head that our ablation isolates.

%======================================================================
\section{Method}\label{sec:method}

\subsection{Task and baselines}\label{sec:task}
Let $\mathcal{D}=\{(s_i,y_i)\}_{i=1}^{N}$ pair each citation context $s_i$
with a label $y_i\in\{\texttt{background},\texttt{method},\texttt{result}\}$.
The baseline is a standard encoder classifier: a pretrained encoder $E$ maps
the tokenised context to a sequence of vectors, the \mbox{[CLS]} position
$h=E(s)[0]$ feeds a linear head,
\begin{equation}\label{eq:baseline}
p(y\mid s)=\mathrm{softmax}\big(W\,h+b\big).
\end{equation}
Three encoders are benchmarked, $E\in\{\texttt{BERT}_{\texttt{base}},
\texttt{SciBERT},\texttt{DeBERTa-v3}_{\texttt{base}}\}$, sharing tokeniser
settings, training schedule and head so that encoder is the only variable.

\subsection{Retrieval-augmented inputs}\label{sec:real}
The augmented variant prepends $k$ evidence contexts retrieved from a pool
$\mathcal{C}$:
\begin{equation}\label{eq:retrieve}
\mathcal{E}_k(s)=\mathrm{top}\text{-}k_{c\in\mathcal{C}}\,
\big[\mathrm{BM25L}(s,c)\big],
\end{equation}
\begin{equation}\label{eq:input}
\tilde{s}=[\,e_1\odot e_2\odot\dots\odot e_k\odot s\,],
\end{equation}
where $\odot$ joins segments with the encoder's separation token. BM25L
scores $s$ against every $c\in\mathcal{C}$ after lowercase, punctuation
stripping and whitespace folding. The parameter $\delta$ (fixed at $0.5$)
lifts the lower bound of term-frequency normalisation, which matters for
citation snippets whose lengths range from a handful of words to sixty.

\subsection{Leakage control}\label{sec:leak}
The pool $\mathcal{C}$ contains training-split contexts only, so nothing is
ever retrieved from the evaluation splits. The query itself is excluded from
its own results through a text-to-index map built at index construction, which
keeps exclusion $O(1)$. A third control, applied for the leak-free ablation,
masks evidence whose label equals the target's, $m(c)=\mathbf{1}[y_c\neq
y_s]$. This does not describe our default protocol---it exists to isolate the
label-copying failure path in one experiment, and it turned out to isolate
an elimination leak instead (Section~\ref{sec:ablation}).

\subsection{Gated dual-pooling head}\label{sec:dual}
An augmented input mixes two kinds of content: evidence segments and the
target sentence. The standard head reads only the pooled
\mbox{[CLS]} position and has no way to tell the two apart. We add a second
view: alongside
\begin{equation}\label{eq:evpool}
h_{\mathrm{cls}}=E(\tilde{s})[0],\qquad
h_{\mathrm{ev}}=\frac{1}{|\mathcal{E}_k(s)|}\sum_{j\in\mathcal{E}_k(s)}
E(\tilde{s})[j],
\end{equation}
where $h_{\mathrm{ev}}$ averages token positions belonging to the evidence
segments, a scalar gate
\begin{equation}\label{eq:gate}
g=\sigma\big(W_g[h_{\mathrm{cls}};h_{\mathrm{ev}}]+b_g\big),\quad
h=g\odot h_{\mathrm{cls}}+(1-g)\odot h_{\mathrm{ev}},
\end{equation}
mixes the two before the linear head of Eq.~\eqref{eq:baseline}. The extra
parameters are one linear layer; nothing else in the model changes. Section
\ref{sec:ablation} shows this head is what keeps the augmented model
trainable at all.

\subsection{Fast, cacheable preprocessing}\label{sec:eff}
Two engineering choices make the pipeline cheap to run and rerun.

\textbf{Vectorised BM25L.} The reference implementation
(\texttt{rank\_bm25}) scores one query term by looping over every corpus
document in Python. We store term frequencies in a sparse matrix and score
each query term by touching only its nonzeros: with
$\mathrm{ctd}=tf/(1-b+b\,|d|/\bar{d})$,
\begin{equation}\label{eq:fast}
\mathrm{score}(s,c)=\sum_{w\in s}\mathrm{idf}(w)\,\frac{tf_{w,c}\,(k_1+1)
\,(\mathrm{ctd}_{w,c}+\delta)}{k_1+\mathrm{ctd}_{w,c}+\delta},
\end{equation}
evaluated as column slices of a CSC matrix. Over all 8{,}243 training queries
the two implementations agree to float64 rounding (max $|\Delta|=5.5\times
10^{-12}$), every query returns the same top-3, and the retrieved evidence is
byte-identical, so the two are interchangeable. The scoring pass drops from
7.0 minutes to 5.5 seconds; the full augmented build, which additionally
assembles the evidence text for all three splits, drops from 8.3 minutes to
11.6 seconds.

\textbf{Disk cache.} Augmented datasets are pickled once per retrieval
configuration (keyed by $k$, leak-free flag and index version). Any rerun or
ablation that reuses the configuration loads from disk in under a second,
and the cached payload is byte-identical to what the fresh pass produces.

%======================================================================
\section{Experiments}\label{sec:experiments}

\subsection{Data and protocol}\label{sec:datasets}
SciCite~\cite{cohan2019} provides 11{,}020 contexts (train 8{,}243, dev 916,
test 1{,}861) over three intents with a skewed distribution:
background $\sim$58\%, method $\sim$29\%, result $\sim$13\%. We report
accuracy and macro-F1; the second matters because the smallest class carries
the most interesting signal. Models select checkpoints by dev macro-F1, three
epochs, AdamW with linear warmup, learning rate $2\times10^{-5}$, weight decay
$0.01$, effective batch 32 (8 per device, 4 accumulation steps), max length
256, seed 42.

The hardware is a single RTX 3060 with 6\,GB of VRAM, which constrains the
recipe: bf16 mixed precision and gradient checkpointing are on; the backbone
loads in fp32 with contiguous parameters. That last detail is load-bearing.
Mixed fp16/fp32 weights, in our runs, converge to majority-class prediction
while the training loss keeps decreasing (53.57\% accuracy, 23.26 macro-F1,
no visible error), so we fix dtype at load time rather than mid-training.

\subsection{Main results}\label{sec:mainresults}

\begin{table}[t]
\caption{SciCite test split. All values are measured, from runs tagged in the
released log. Retrieval uses $k{=}3$; ``dual'' marks the gated dual-pooling
head.}
\label{tab:main}
\begin{center}
\small
\begin{tabular}{lcc}
\toprule
\textbf{Model} & \textbf{Acc\,(\%)} & \textbf{Macro-F1\,(\%)}\\
\midrule
BERT-base                    & 85.71 & 84.11\\
SciBERT                      & 86.30 & 85.11\\
DeBERTa-v3-base              & \textbf{86.89} & \textbf{85.49}\\
\midrule
DeBERTa + retrieval (CLS)    & 59.97 & 38.73\\
DeBERTa + retrieval (dual, run 1) & 84.95 & 83.30\\
DeBERTa + retrieval (dual, run 2) & 86.14 & 84.71\\
\bottomrule
\end{tabular}
\end{center}
\end{table}

Table~\ref{tab:main} holds the core comparison, and the reading is
uneventful until the last two rows.

The encoder ordering is as expected (DeBERTa-v3-base ahead of SciBERT ahead
of BERT) and the margins are small but consistent. A strong backbone sets a
high bar for augmentation; two of three baselines clear 86\% accuracy.

Retrieval with a bare CLS head fails outright: 59.97\% accuracy and 38.73
macro-F1 is barely above the majority-class rate on this split. The pooled
[CLS] vector has to summarise a sequence in which three retrieved sentences
precede the target, and it does not. Training completes without complaint,
which makes the failure easy to miss.

The dual head repairs this. Both dual-pooling runs land within 0.8 points of
the DeBERTa baseline (86.14\% / 84.71\% and 84.95\% / 83.30\%, two runs of the
same configuration, shown separately to indicate run-to-run variance). The
repaired augmented model matches, and does not beat, its backbone. Under the
leakage controls of Section~\ref{sec:leak}, we found no configuration with
label-blind inputs in which retrieval improved either metric over the
corresponding plain encoder; the one label-conditioned exception is treated
in Section~\ref{sec:ablation} as a leak in its own right.

\subsection{Head ablation and stability}\label{sec:ablation}

\begin{table}[t]
\caption{Head, precision and leakage ablation, DeBERTa-v3 with retrieval
$k{=}3$. Same seed, same data, one variable at a time. The leak-free row is
an elimination leak, not a gain.}
\label{tab:head}
\begin{center}
\small
\begin{tabular}{lcc}
\toprule
\textbf{Configuration} & \textbf{Acc\,(\%)} & \textbf{Macro-F1\,(\%)}\\
\midrule
CLS head                    & 59.97 & 38.73\\
Dual head (ours)            & \textbf{86.14} & \textbf{84.71}\\
\midrule
fp16/fp32 mix, no fix       & 53.57 & 23.26\\
fp32 load, contiguous fix   & \textbf{86.14} & \textbf{84.71}\\
\midrule
Leak-free mask (dual head)  & 96.40 & 95.04\\
\bottomrule
\end{tabular}
\end{center}
\end{table}

Table~\ref{tab:head} separates the two fixes our pipeline needed. The
dual-pooling head accounts for 26 points of accuracy against the CLS baseline
on identical data. The dtype fix accounts for a further silent failure: with
half-precision backbone weights and an fp32 head, the same run collapses
toward the majority class while its loss curve looks healthy. We flag both
here because each failure is invisible from the loss alone, and both are
cheap to prevent once known.

The leak-free row posts the best numbers in the paper and is the least
trustworthy of them. The mask of Section~\ref{sec:leak} removes every
candidate whose label equals the target's, so the $k{=}3$ neighbours that
survive carry only the other two classes; the target label is then the
class absent from the input, and the encoder learns the elimination rule
(96.40\% test accuracy, 97.88\% train accuracy---the shortcut transfers
perfectly, which is what a shortcut looks like). The lesson is not that
retrieval finally helped. It is that a leakage control conditioned on the
label changes the input distribution in label-correlated ways: anything
conditioned on $y$ at input-construction time is a channel, whether it adds
evidence or removes it. Our default protocol retrieves label-blind, so the
honest comparison is rows 1--2.

\subsection{Efficiency}\label{sec:efficiency}

\begin{table}[t]
\caption{Measured cost on the RTX 3060 that produced all reported results.
Top: retrieval preprocessing, reference \texttt{rank\_bm25} vs.\ our sparse
scorer. Bottom: the augmented model against its plain backbone.}
\label{tab:efficiency}
\begin{center}
\small
\begin{tabular}{lcc}
\toprule
\textbf{Metric} & \textbf{Reference / plain} & \textbf{Ours / augmented}\\
\midrule
Scoring pass, 8{,}243 queries & 7.0\,min   & 5.5\,s\\
Full augmented build          & 8.3\,min   & 11.6\,s\\
Cached re-build              & ---        & 0.5\,s\\
\midrule
Train throughput (samp/s)     & 34.1       & 34.4\\
Eval throughput (samp/s)      & 157.5      & 153.5\\
\bottomrule
\end{tabular}
\end{center}
\end{table}

The preprocessing numbers come from a benchmark that scores every training
context against the pool with both implementations and compares results
(Section~\ref{sec:eff}). Scoring agrees to float64 rounding, every query
returns the same top-3, and the retrieved evidence is byte-identical for all
8{,}243 queries. The full augmented build drops from 8.3 minutes to 11.6
seconds; the disk cache then removes even that from every rerun which
reuses the configuration (0.5\,s to reload).

The training comparison needs one caveat to read fairly. The plain baseline
pads every context to 256 tokens although the median context is 48 tokens
long, so it pushes 2.1M tokens through the encoder per epoch, most of it
padding; the augmented sequences fill the same budget (median 256), for the
same 2.1M tokens per epoch. On that shared token budget the two models train
at the same throughput (34.1 vs.\ 34.4 samples/s, measured in a controlled
benchmark with the exact configuration of Section~\ref{sec:datasets}) and
evaluate within 2.5\% of each other (157.5 vs.\ 153.5 samples/s). For
practical purposes the augmentation is free at training time on this
hardware; it costs 12 seconds before training starts, once per
configuration.

The consequence for experimental practice is straightforward. The full grid
behind this paper, three baselines, three augmented runs and the diagnostics
behind each ablation, was computed on one 6\,GB GPU in a working session,
and any of it can be reproduced from a cold start in minutes rather than
hours.

\subsection{Why does retrieval not help here?}\label{sec:discussion}
The evidence BM25L returns is lexical. On citation contexts, lexical
similarity tracks label-specific phrasing (``compared with'',
``following the approach of''), which the encoder already reads off the
target sentence directly. The retrieved neighbours add, at best, a redundant
restatement of what the backbone can extract in one pass. What they add at
worst is confusion, which is exactly what the CLS-head collapse shows. The
dual head recovers the baseline rather than exceeding it because the gate can
learn to down-weight the evidence entirely; nothing in the augmented input
contains information the plain input lacks.

This is a claim about CIC under a strong in-domain encoder, not about
retrieval augmentation in general. QA needs the corpus because the question
does not contain the answer; here the sentence, for most instances, does
contain the label signal. Where that stops holding (multi-intent
taxonomies, cross-domain transfer, contexts whose meaning depends on the
cited abstract), the balance may shift. Our protocol, with its leakage
controls and head ablation, is the reusable part: it is how one finds out.

%======================================================================
\section{Conclusion}\label{sec:conclusion}

We set out to test whether BM25L retrieval augmentation improves citation
intent classification under leakage controls, with strong pretrained
backbones, on public data. It does not: DeBERTa-v3-base at 86.89\% /
85.49\% outperforms every augmented variant we measured, and the single
variant that outscores it does so through an elimination leak introduced by
its own label-conditioned mask. The study's value
lies in what it isolates. A standard CLS head silently breaks on augmented
inputs, and a gated dual-pooling head recovers them, a 26-point swing that
anyone running this recipe should know about before drawing conclusions from
a failed or successful run. And the cost of finding such things out is now
small: vectorised scoring plus caching brings a full preprocessing pass
under 12 seconds and a cached rebuild under half a second, with training
throughput unchanged on 6\,GB of VRAM.

Dense retrievers, cited-abstract fusion, and harder taxonomies such as
MultiCite are the natural next tests for the protocol.

%======================================================================
\section*{Data and code availability}
SciCite is distributed under Apache-2.0 by AllenAI
(\url{https://github.com/allenai/SciCite}). Our implementation,
configurations and the raw metrics behind every number in this paper are
released as the \emph{SciCite-REAL} repository, Apache-2.0, at
\url{https://github.com/Merthyr04/SciCite-REAL}.

%======================================================================
\begin{thebibliography}{99}

\bibitem{cohan2019} A.~Cohan, W.~Ammar, M.~van~Zuylen, and F.~Cady,
``Structural scaffolds for citation intent classification in scientific
publications,'' in \emph{Proc.\ NAACL}, 2019.

\bibitem{beltagy2019} I.~Beltagy, K.~Lo, and A.~Cohan, ``SciBERT: A
pretrained language model for scientific text,'' in \emph{Proc.\ EMNLP},
2019.

\bibitem{he2021} P.~He, J.~Gao, and W.~Chen, ``DeBERTaV3: Improving DeBERTa
using ELECTRA-style pre-training with gradient-disentangled embedding
sharing,'' arXiv:2111.09543, 2021.

\bibitem{devlin2019} J.~Devlin, M.-W.~Chang, K.~Lee, and K.~Toutanova,
``BERT: Pre-training of deep bidirectional transformers for language
understanding,'' in \emph{Proc.\ NAACL}, 2019.

\bibitem{lewis2020} P.~Lewis, E.~Perez, A.~Piktus, et~al.,
``Retrieval-augmented generation for knowledge-intensive NLP tasks,'' in
\emph{Proc.\ NeurIPS}, 2020.

\bibitem{karpukhin2020} V.~Karpukhin, B.~O\u{g}uz, S.~Min, et~al., ``Dense
passage retrieval for open-domain question answering,'' in \emph{Proc.\
EMNLP}, 2020.

\bibitem{guu2020} K.~Guu, K.~Lee, Z.~Tung, P.~Pasupat, and M.-W.~Chang,
``REALM: Retrieval-augmented language model pre-training,'' in \emph{Proc.\
ICML}, 2020.

\bibitem{robertson2009} S.~Robertson and H.~Zaragoza, ``The probabilistic
relevance framework: BM25 and beyond,'' \emph{Foundations and Trends in
Information Retrieval}, vol.~3, no.~4, pp.~333--389, 2009.

\bibitem{lv2011} Y.~Lv and C.~Zhai, ``Lower-bounding term frequency
normalization,'' in \emph{Proc.\ CIKM}, 2011.

\bibitem{jurgens2018} D.~Jurgens, S.~Kumar, R.~Hoover, D.~McFarland, and
D.~Jurafsky, ``Measuring the evolution of a scientific field through
citation Frames,'' \emph{Transactions of the ACL}, vol.~6, 2018.

\bibitem{teufel2006} S.~Teufel, A.~Siddharthan, and D.~Tidhar, ``Automatic
classification of citation function,'' in \emph{Proc.\ EMNLP}, 2006.

\bibitem{peters2018} M.~Peters, M.~Neumann, M.~Iyyer, et~al., ``Deep
contextualized word representations,'' in \emph{Proc.\ NAACL}, 2018.

\bibitem{clark2020} K.~Clark, M.-T.~Luong, Q.~V.~Le, and C.~D.~Manning,
``ELECTRA: Pre-training text encoders as discriminators rather than
generators,'' in \emph{Proc.\ ICLR}, 2020.

\end{thebibliography}

\end{document}